\documentclass[a4paper,10pt]{article}
\usepackage{fontspec}
\setmainfont{Libertinus Serif}

\usepackage[a4paper,margin=1.2cm,top=1.0cm,bottom=1.0cm]{geometry}
\usepackage[ngerman]{babel}
\usepackage{amsmath,amssymb}
\usepackage{xcolor}
\usepackage{tikz}
\usepackage{fontawesome5}
\usepackage[skins]{tcolorbox}
\usepackage{enumitem}

\definecolor{primaryBlue}{HTML}{143264}
\definecolor{lvlGreen}{RGB}{25, 125, 55}
\definecolor{lvlBlue}{RGB}{20, 80, 160}
\definecolor{lvlPurple}{RGB}{115, 30, 135}

\definecolor{step1Bg}{HTML}{FFFBEB}
\definecolor{step1Border}{HTML}{FDE68A}
\definecolor{step1Text}{HTML}{92400E}

\definecolor{step2Bg}{HTML}{EFF6FF}
\definecolor{step2Border}{HTML}{BFDBFE}
\definecolor{step2Text}{HTML}{1E40AF}

\definecolor{step3Bg}{HTML}{ECFDF5}
\definecolor{step3Border}{HTML}{A7F3D0}
\definecolor{step3Text}{HTML}{065F46}

\pagestyle{empty}
\setlength{\parindent}{0pt}

\begin{document}

% ==============================================================================
% SEITE 1: TIPP-KARTEN PROZENTRECHNUNG (S. 223 - 224)
% ==============================================================================

\noindent
{\color{primaryBlue}\bfseries\normalsize \faBook\enskip Gestufte Hilfen (Printversion): Prozentrechnung}
\hfill
{\footnotesize\color{primaryBlue!80} \textbf{Mathematik 6D} \quad$|$\quad Herr Dayi}
\vspace{1.5mm}
{\color{primaryBlue!40}\hrule height 0.6pt}
\vspace{2.5mm}

% --- AUFGABE S. 223 NR. 2 ---
\begin{tcolorbox}[enhanced, colback=lvlGreen!5, colframe=lvlGreen!50, arc=2mm, boxrule=0.8pt, left=3mm, right=3mm, top=2mm, bottom=2mm]
    {\bfseries\color{lvlGreen} \faStar\enskip S. 223, Nr. 2 (Basis): Brüche in Prozent umwandeln}\\[1mm]
    {\footnotesize\textit{Wandle die Brüche $\frac{1}{2}, \frac{1}{5}, \frac{3}{5}, \frac{1}{4}, \frac{2}{25}, \dots$ in Prozent um.}}
\end{tcolorbox}

\vspace{1.5mm}

\begin{tcolorbox}[colback=step1Bg, colframe=step1Border, arc=1.8mm, boxrule=0.7pt, left=3mm, right=3mm, top=1.8mm, bottom=1.8mm]
    {\color{step1Text}\small\bfseries\faEye\enskip Stufe 1: Anschauung}\\[1mm]
    {\footnotesize Prozent bedeutet \textbf{„von Hundert“}. Du musst den Bruch so erweitern, dass im Nenner $\mathbf{100}$ steht. Dann liest du die Prozentzahl direkt im Zähler ab!}
\end{tcolorbox}

\begin{tcolorbox}[colback=step2Bg, colframe=step2Border, arc=1.8mm, boxrule=0.7pt, left=3mm, right=3mm, top=1.8mm, bottom=1.8mm]
    {\color{step2Text}\small\bfseries\faLightbulb\enskip Stufe 2: Denkanstoß}\\[1mm]
    {\footnotesize Mit welcher Zahl musst du den Nenner multiplizieren, um 100 zu erhalten?\\
    Nenner 2: $\cdot 50$ \quad$|$\quad Nenner 5: $\cdot 20$ \quad$|$\quad Nenner 4: $\cdot 25$ \quad$|$\quad Nenner 25: $\cdot 4$ \quad$|$\quad Nenner 50: $\cdot 2$.}
\end{tcolorbox}

\begin{tcolorbox}[colback=step3Bg, colframe=step3Border, arc=1.8mm, boxrule=0.7pt, left=3mm, right=3mm, top=1.8mm, bottom=1.8mm]
    {\color{step3Text}\small\bfseries\faTh\enskip Stufe 3: Rechengerüst \& Kontrolltipp}\\[1mm]
    {\footnotesize
    Nutze das Erweiterungsgerüst für dein Heft:\\
    $\frac{\text{Zähler}}{\text{Nenner}} = \frac{\text{Zähler} \cdot \text{Erweiterungszahl}}{\text{Nenner} \cdot \text{Erweiterungszahl}} = \frac{\Box}{100} = \Box\,\%$.\\[1mm]
    \textit{Kontrolltipp:} Achte darauf, dass im Nenner immer genau $100$ steht. Dann liest du die Prozentzahl direkt oben im Zähler ab!}
\end{tcolorbox}

\vspace{3.5mm}

% --- AUFGABE S. 224 NR. 9 ---
\begin{tcolorbox}[enhanced, colback=lvlBlue!5, colframe=lvlBlue!50, arc=2mm, boxrule=0.8pt, left=3mm, right=3mm, top=2mm, bottom=2mm]
    {\bfseries\color{lvlBlue} \faStar\faStar\enskip S. 224, Nr. 9 (Standard): Welche Angabe passt nicht?}\\[1mm]
    {\footnotesize\textit{Vergleiche: a) jeder Fünfte; $20\,\%$; $3$ von $15$; $\frac{1}{4}$.}}
\end{tcolorbox}

\vspace{1.5mm}

\begin{tcolorbox}[colback=step1Bg, colframe=step1Border, arc=1.8mm, boxrule=0.7pt, left=3mm, right=3mm, top=1.8mm, bottom=1.8mm]
    {\color{step1Text}\small\bfseries\faEye\enskip Stufe 1: Anschauung}\\[1mm]
    {\footnotesize Wandle alle vier Ausdrücke in dieselbe Darstellungsform um (am einfachsten alle in Prozent!).}
\end{tcolorbox}

\begin{tcolorbox}[colback=step2Bg, colframe=step2Border, arc=1.8mm, boxrule=0.7pt, left=3mm, right=3mm, top=1.8mm, bottom=1.8mm]
    {\color{step2Text}\small\bfseries\faLightbulb\enskip Stufe 2: Denkanstoß}\\[1mm]
    {\footnotesize Kürze $\frac{3}{15}$ zuerst mit 3! Welcher Stammbruch entsteht? Erweitere diesen dann auf Nenner 100.}
\end{tcolorbox}

\begin{tcolorbox}[colback=step3Bg, colframe=step3Border, arc=1.8mm, boxrule=0.7pt, left=3mm, right=3mm, top=1.8mm, bottom=1.8mm]
    {\color{step3Text}\small\bfseries\faCheckSquare\enskip Stufe 3: Rechengerüst \& Kontrolltipp}\\[1mm]
    {\footnotesize
    • „Jeder Fünfte“ als Bruch: $\frac{1}{5} = \frac{\dots}{100} = \dots\,\%$.\\
    • „3 von 15“ $\xrightarrow{\text{mit 3 kürzen}} \frac{1}{\dots} = \frac{\dots}{100} = \dots\,\%$.\\
    • $\frac{1}{4} = \frac{\dots}{100} = \dots\,\%$.\\[1mm]
    \textit{Kontrolltipp:} Genau drei Angaben ergeben denselben Prozentwert ($20\,\%$). Die abweichende Angabe passt nicht!}
\end{tcolorbox}

\newpage

% ==============================================================================
% SEITE 2: TIPP-KARTEN BRÜCHE ALS QUOTIENT (S. 227)
% ==============================================================================

\noindent
{\color{primaryBlue}\bfseries\normalsize \faDivide\enskip Gestufte Hilfen (Printversion): Brüche als Quotient}
\hfill
{\footnotesize\color{primaryBlue!80} \textbf{Mathematik 6D} \quad$|$\quad Herr Dayi}
\vspace{1.5mm}
{\color{primaryBlue!40}\hrule height 0.6pt}
\vspace{2.5mm}

% --- AUFGABE S. 227 NR. 1 ---
\begin{tcolorbox}[enhanced, colback=lvlGreen!5, colframe=lvlGreen!50, arc=2mm, boxrule=0.8pt, left=3mm, right=3mm, top=2mm, bottom=2mm]
    {\bfseries\color{lvlGreen} \faStar\enskip S. 227, Nr. 1 (Basis): Division als Bruch schreiben}\\[1mm]
    {\footnotesize\textit{Schreibe als vollständig gekürzten Bruch: a) $5 : 2$, b) $2 : 7$, c) $15 : 5$, i) $18 : 12$.}}
\end{tcolorbox}

\vspace{1.5mm}

\begin{tcolorbox}[colback=step1Bg, colframe=step1Border, arc=1.8mm, boxrule=0.7pt, left=3mm, right=3mm, top=1.8mm, bottom=1.8mm]
    {\color{step1Text}\small\bfseries\faEye\enskip Stufe 1: Anschauung}\\[1mm]
    {\footnotesize Die Grundregel lautet: Erste Zahl nach \textbf{oben} in den Zähler, zweite Zahl nach \textbf{unten} in den Nenner (z.\,B. $3 : 4 = \frac{3}{4}$).}
\end{tcolorbox}

\begin{tcolorbox}[colback=step2Bg, colframe=step2Border, arc=1.8mm, boxrule=0.7pt, left=3mm, right=3mm, top=1.8mm, bottom=1.8mm]
    {\color{step2Text}\small\bfseries\faLightbulb\enskip Stufe 2: Denkanstoß}\\[1mm]
    {\footnotesize Prüfe nach dem Aufschreiben: Geht die Division ohne Rest auf (z.\,B. $15:5 = \frac{15}{5} = 3$)? Kannst du Zähler und Nenner noch kürzen?}
\end{tcolorbox}

\begin{tcolorbox}[colback=step3Bg, colframe=step3Border, arc=1.8mm, boxrule=0.7pt, left=3mm, right=3mm, top=1.8mm, bottom=1.8mm]
    {\color{step3Text}\small\bfseries\faTh\enskip Stufe 3: Rechengerüst \& Kontrolltipp}\\[1mm]
    {\footnotesize
    $\text{Zahl 1} : \text{Zahl 2} = \frac{\text{Zahl 1}}{\text{Zahl 2}}$. Bei unechten Brüchen (Zähler $>$ Nenner) als gemischte Zahl schreiben.\\
    $\bullet$ Bei $18 : 12 = \frac{18}{12}$: Kürze Zähler und Nenner mit $6$!\\
    \textit{Kontrolltipp:} Genau eine der drei Basis-Aufgaben ergibt eine natürliche Zahl ohne Rest!}
\end{tcolorbox}

\vspace{3.5mm}

% --- AUFGABE S. 227 NR. 2 ---
\begin{tcolorbox}[enhanced, colback=lvlBlue!5, colframe=lvlBlue!50, arc=2mm, boxrule=0.8pt, left=3mm, right=3mm, top=2mm, bottom=2mm]
    {\bfseries\color{lvlBlue} \faStar\faStar\enskip S. 227, Nr. 2 (Standard): Kistensortierung}\\[1mm]
    {\footnotesize\textit{Sortiere die Kärtchen in die Kisten $\frac{2}{3}$, $\frac{3}{2}$ und $\frac{3}{4}$ ein: $2:3$, $3:2$, $9:12$, $12:16$, $8:12$, $40:60$, $18:12$, $48:32$, $100:150$.}}
\end{tcolorbox}

\vspace{1.5mm}

\begin{tcolorbox}[colback=step1Bg, colframe=step1Border, arc=1.8mm, boxrule=0.7pt, left=3mm, right=3mm, top=1.8mm, bottom=1.8mm]
    {\color{step1Text}\small\bfseries\faEye\enskip Stufe 1: Anschauung}\\[1mm]
    {\footnotesize Schreibe jedes Kärtchen als Bruch auf und kürze ihn vollständig!}
\end{tcolorbox}

\begin{tcolorbox}[colback=step2Bg, colframe=step2Border, arc=1.8mm, boxrule=0.7pt, left=3mm, right=3mm, top=1.8mm, bottom=1.8mm]
    {\color{step2Text}\small\bfseries\faLightbulb\enskip Stufe 2: Denkanstoß}\\[1mm]
    {\footnotesize Schneller Check für Kiste $\frac{3}{2}$: Hier ist der Zähler größer als der Nenner (Wert $>1$). Alle Kärtchen mit erster Zahl $>$ zweiter Zahl gehören hierhin!}
\end{tcolorbox}

\begin{tcolorbox}[colback=step3Bg, colframe=step3Border, arc=1.8mm, boxrule=0.7pt, left=3mm, right=3mm, top=1.8mm, bottom=1.8mm]
    {\color{step3Text}\small\bfseries\faCheckSquare\enskip Stufe 3: Kontrollhilfe zur Einordnung}\\[1mm]
    {\footnotesize
    Kürze jeden Quotienten mit dem gemeinsamen Teiler auf die Grundform:\\
    z.\,B. $8:12 = \frac{8}{12} \xrightarrow{\text{mit 4 kürzen}} \frac{2}{3}$ \quad$|$\quad $48:32 = \frac{48}{32} \xrightarrow{\text{mit 16 kürzen}} \frac{3}{2}$.\\[1mm]
    \textit{Kontrollhilfe:} In Kiste $\frac{3}{2}$ gehören genau \textbf{3 Kärtchen}, in Kiste $\frac{3}{4}$ genau \textbf{4 Kärtchen} und in Kiste $\frac{2}{3}$ genau \textbf{5 Kärtchen}!}
\end{tcolorbox}

\vspace{3.5mm}

% --- AUFGABE S. 228 NR. 7 ---
\begin{tcolorbox}[enhanced, colback=lvlPurple!5, colframe=lvlPurple!50, arc=2mm, boxrule=0.8pt, left=3mm, right=3mm, top=2mm, bottom=2mm]
    {\bfseries\color{lvlPurple} \faStar\faStar\faStar\enskip S. 228, Nr. 7 (Extra): Finde den Fehler!}\\[1mm]
    {\footnotesize\textit{Finde und korrigiere den Fehler: a) $8 : 5 = \frac{5}{8}$, b) Mädchen zu Jungen 3:6 $\to$ die Hälfte Mädchen, c) Quotient aus 8 und 8 ist $\frac{8}{8} = 0$.}}
\end{tcolorbox}

\vspace{1.5mm}

\begin{tcolorbox}[colback=step1Bg, colframe=step1Border, arc=1.8mm, boxrule=0.7pt, left=3mm, right=3mm, top=1.8mm, bottom=1.8mm]
    {\color{step1Text}\small\bfseries\faEye\enskip Stufe 1: Anschauung}\\[1mm]
    {\footnotesize Prüfe: Welche Zahl steht vorne (Zähler) und welche hinten (Nenner)? Wie viele Kinder sind es insgesamt in Klasse 5b? Was ist eine Zahl geteilt durch sich selbst?}
\end{tcolorbox}

\begin{tcolorbox}[colback=step2Bg, colframe=step2Border, arc=1.8mm, boxrule=0.7pt, left=3mm, right=3mm, top=1.8mm, bottom=1.8mm]
    {\color{step2Text}\small\bfseries\faLightbulb\enskip Stufe 2: Denkanstoß}\\[1mm]
    {\footnotesize a) Zähler und Nenner wurden vertauscht! b) Insgesamt gibt es $3 + 6 = 9$ Kinder (Anteil ist $\frac{3}{9}$)! c) $8:8 = 1$ Ganzes, nicht $0$!}
\end{tcolorbox}

\begin{tcolorbox}[colback=step3Bg, colframe=step3Border, arc=1.8mm, boxrule=0.7pt, left=3mm, right=3mm, top=1.8mm, bottom=1.8mm]
    {\color{step3Text}\small\bfseries\faPenNib\enskip Stufe 3: Formulierungshilfe zur Begründung}\\[1mm]
    {\footnotesize
    Formuliere deine Erklärung mit den mathematischen Fachbegriffen so:\\
    • zu a): „Der Fehler liegt in der Reihenfolge: Die erste Zahl gehört in den \dots, die zweite Zahl in den \dots“\\
    • zu b): „Der Fehler ist, dass die Klasse insgesamt $3 + 6 = \dots$ Kinder hat. Der Anteil ist also $\frac{3}{\dots}$ und nicht die Hälfte.“\\
    • zu c): „Eine Zahl dividiert durch sich selbst ergibt immer \dots Ganze(s).“}
\end{tcolorbox}

\end{document}
